% Solutions to the exercises of Chapter 29.
\begin{solutions}
  \item Draw a random state (IMU, one clone, one feature). Apply $p \mapsto R_zp + t$, $v \mapsto R_zv$, $R \mapsto R_zR$ (and the same to the clone and the feature) with $R_z = \Exp(\alpha e_z)$,
        leaving the biases unchanged. Compute $x'\bminus x$ with the local attitude difference $\Log(\hat R\T R')$, and compare it with $N[t;\ \alpha]$. With $\|t\|, |\alpha| \approx 10^{-3}$, the
        discrepancy is $4\times10^{-5}$ against entries of $3\times10^{-2}$. With $10^{-5}$ it is $2\times10^{-10}$ against $6\times10^{-5}$: it falls quadratically, so $N$ is the exact first-order map. The
        attitude rows come out as $\hat R\T e_z$ and $\hat R_i\T e_z$. With a global attitude error they would be the constant $e_z$, a reminder that $N$ must be written in the
        filter's own error coordinates.

  \item For \emph{every} state, translating all positions and features leaves both the dynamics and the measurements unchanged. The nominal kinematics never use
        the position itself, only its rate, and the camera depends on positions only through differences $p_f - p_{wb_i}$. Differentiating these identities
        at \emph{any} linearization point gives $\Phi_kN_t = N_t$ and $H_kN_t = 0$, or $H_{x,j}N_t + H_{f,j} = 0$ and hence $H_{o,j}N_t = 0$ after projection. The translation columns
        $N_t = (I, 0, 0, 0, 0, I, 0, \dots)$ contain no estimates at all, so the product telescopes trivially, whatever estimates the individual factors used:
        \[
          H_k\Phi_{k-1}\cdots\Phi_1N_t = H_kN_t = 0 .
        \]
        The yaw column contains $\hat p$, $\hat v$, $\hat R$ and the clone estimates. The $N$ that one factor maps into is then not the $N$ that the next factor annihilates, unless the
        estimates agree. (The pose measurement of Milestones 6--7 is different: it observes absolute position, so translations are observable there.)

  \item The state is $x = (p, \psi, \ell_1, \ell_2) \in \R^7$ and the output $h = (R\T(\ell_1 - p), R\T(\ell_2 - p)) \in \R^4$, with
        $\nabla h_i = \big[-R\T\ \ -\Omega h_i\ \ \cdots\ \ R\T\ \ \cdots\big]$, the $R\T$ in the columns of $\ell_i$. As before, $L_{f_v}h_i = -e_1$ is constant and $L_{f_\omega}h_i = -\Omega h_i$, whose gradient
        $-\Omega\nabla h_i$ lies in the row space of $\nabla h_i$; the same holds for all higher derivatives. So $d\mathcal{O}$ is the row space of $[\nabla h_1;\ \nabla h_2]$. Its landmark
        columns form the invertible $\blkdiag(R\T, R\T)$, so the rank is 4. (Numerically, with the closed-form Lie derivatives, the singular values are $12.2$, $3.0$, $2.3$,
        $1.7$, and then exactly zero three times.) Three directions are unobservable, the same three as with one landmark:
        \[
          N_t = (I_2,\ 0,\ I_2,\ I_2), \qquad N_r = (\Omega p,\ 1,\ \Omega\ell_1,\ \Omega\ell_2),
        \]
        since $\nabla h_iN_r = -R\T\Omega p - \Omega R\T(\ell_i - p) + R\T\Omega\ell_i = 0$ by $R\T\Omega = \Omega R\T$. Each new landmark adds two states and two independent outputs, so landmarks never reduce
        the number of unobservable directions. Only known landmarks do.

  \item The measurement depends on the state only through $d = R(\psi)\T(\ell_j - p)$, so by the chain rule $H_k^{\mathrm{rb}} = J_z(\hat d_k)\,H_k$, with $H_k$ the Cartesian Jacobian
        \eqref{eq:slam-H} and
        \[
          J_z = \frac{\partial(\|d\|, \atantwo(d_y, d_x))}{\partial d} = \begin{bmatrix} d\T/\rho \\ (\Omega d)\T/\rho^2 \end{bmatrix}, \qquad \rho = \|d\|, \qquad \det J_z = \frac1\rho \ne 0 .
        \]
        Multiplying \cref{thm:slam-obs} on the left by $J_z$:
        \[
          H_k^{\mathrm{rb}}\Phi_{k-1}\cdots\Phi_1N_t = 0, \qquad
          H_k^{\mathrm{rb}}\Phi_{k-1}\cdots\Phi_1N_r = J_z(\hat d_k)\,\hat R_k\T\Omega\big(D_k + \bar\ell_j - \hat\ell_j\big).
        \]
        The leak is the Cartesian one premultiplied by an invertible matrix, so it vanishes under exactly the same conditions, and FEJ cures it in the same way. This
        holds for any measurement of the body-frame landmark position: only the factor $J_z$ changes.

  \item Propagate a copy of $N$ alongside $P$: multiply it by each $\Phi_k$ the filter uses and append $J_xN$ when a landmark is initialized (the filter's own
        augmentation Jacobian). At every update, print $\|HN\|$ for the Jacobian $H$ the filter actually uses. Over five runs of the \cref{fig:fej} scenario:
        \begin{center}\small
        \begin{tabular}{@{}l c c@{}}
        \toprule
        filter & $\max\|HN_t\|$ & $\|HN_r\|$ (median / max) \\
        \midrule
        standard EKF & 0 & 0.17 / 0.9 \\
        FEJ-EKF & 0 & $9\times10^{-16}$ / $4\times10^{-15}$ \\
        ideal EKF & 0 & $9\times10^{-16}$ / $4\times10^{-15}$ \\
        \bottomrule
        \end{tabular}
        \end{center}
        Only the standard EKF violates the condition, and only in the rotation direction, by an amount that tracks the accumulated corrections. One subtlety
        appeared on the way. If the ideal filter evaluates the landmark-initialization Jacobian with the \emph{noisy} measurement, $\Omega R(\psi)z$ instead of $\Omega(\ell_j - p)$, the landmark
        enters $N$ at $\ell_j + Rn$ while $H$ uses $\ell_j$. Then $\|HN_r\| \approx 0.12$, the size of the measurement noise, and the ``ideal'' filter gains a little heading information (its
        final RMS heading error drops from $2.8^\circ$ to $2.5^\circ$). Every Jacobian that involves a state must use the same linearization point, the augmentation Jacobian
        included.

  \item Each row is separate: minimize $\|d\|^2$ subject to $dN = c$, with $d \in \R^{1\times n}$ and $N \in \R^{n\times m}$ of full column rank. Split $d = d_\parallel + d_\perp$ with $d_\parallel = yN\T$ in the row space of
        $N\T$ and $d_\perp N = 0$. The constraint sees only $d_\parallel$: $yN\T N = c$, so $y = c(N\T N)\inv$ and $d_\parallel = c(N\T N)\inv N\T$ is the same for every feasible $d$. Since $d_\parallel d_\perp\T = yN\T d_\perp\T = 0$,
        \[
          \|d\|^2 = \|d_\parallel\|^2 + \|d_\perp\|^2 \ge \|d_\parallel\|^2,
        \]
        with equality exactly when $d_\perp = 0$. The unique minimizer is $d = c(N\T N)\inv N\T$, which with $c = -hN$ (for $H$) or $c = n_{k+1} - hN_k$ (for $\Phi$) gives \cref{prop:oc-ekf}.

  \item At rest the velocity is zero, so use the pseudo-measurement $z = 0 = v_{wb} + n_{\mathrm{zv}}$ with $R_{\mathrm{zv}} = \sigma_{\mathrm{zv}}^2I_3$. The residual is $r = 0 - \hat v$, and since the velocity error is additive
        ($v = \hat v + \delta v$),
        \[
          H = \begin{bmatrix} 0_3 & I_3 & 0_3 & 0_3 & 0_3 & 0_{3\times6N} \end{bmatrix}.
        \]
        A companion \emph{zero-rate} update uses $\omega_m = b_g + n_g$ at rest: a direct measurement of the gyro bias, with $H$ selecting the $\delta b_g$ columns. Safety rules:
        \begin{itemize}
          \item apply it only when stillness is certain: accelerometer and gyro variance over a window below thresholds, and negligible feature motion.
                A hovering drone is not stationary, and a false zero-velocity update is a confident wrong measurement;
          \item keep $\sigma_{\mathrm{zv}}$ realistic (centimeters per second, to cover vibration);
          \item apply it at a modest rate, not at every IMU sample. All these updates rest on one stillness hypothesis, so their errors are not
                independent, and 200 of them per second would make the filter overconfident.
        \end{itemize}
        For EuRoC it applies only before take-off. It bounds velocity and position drift at rest, but adds nothing about yaw.

  \item Following~\cite{li2013msckf2} and \cref{sec:fej}, three kinds of Jacobian need fixed linearization points:
        \begin{enumerate}[label=(\alph*)]
          \item the IMU transition matrix from $t_k$ to $t_{k+1}$, evaluated with the propagated estimates at both ends, $\hat x_{k|k-1}$ and $\hat x_{k+1|k}$, so that consecutive matrices telescope.
                \ifsolutions This includes the attitude block, which must be $\hat R_{k+1|k}\T\hat R_{k|k-1}$ instead of $\Exp(-\hat\omega\Dt)$, so that it maps $\hat R_{k|k-1}\T e_z$ to $\hat R_{k+1|k}\T e_z$;
                \else This includes the attitude block, which must map $\hat R_{k|k-1}\T e_z$ to $\hat R_{k+1|k}\T e_z$ (\cref{sec:fej});\fi
          \item the camera Jacobians with respect to each clone, evaluated at the clone's \emph{first estimate}, its pose when it was cloned (the propagated estimate at the image
                time), while the residual uses the current estimate;
          \item the feature Jacobians, evaluated at one estimate of the feature for all of its observations (any single estimate works, because the feature is
                eliminated).
        \end{enumerate}
        Storage without allocation: give the clone record a second pose, e.g.\ \code{struct Clone \{ nanoseconds t; Vector3d p, p\_fej; Quaterniond q, q\_fej; uint64\_t id; \}} in the
        fixed-capacity ring buffer of \code{kMaxClones} entries. The FEJ fields are written once, at cloning, and skipped by injection. Marginalization drops them with the
        clone. For the IMU, keep one member copy of the previous propagated state $(\hat p, \hat v, \hat q)_{k|k-1}$: compute $\Phi_k$ from it and the newly propagated state, then
        overwrite it with the newly propagated state \emph{before} any update is applied. All of it is fixed-size, so E4 is unaffected. The test is the one of the
        implementation note: $\|H_kN_k\|$ and $\|\Phi_kN_k - N_{k+1}\|$ at rounding level on a synthetic run.
\end{solutions}
